\section{Environnement du stage}

    \begin{verbatim}
- Stage à lieu au LIFO, dans l'équipe CA, encadré par ...
- Stage de fin de master d'informatique Minerve GPEx
- Challenge@EvalLLM comme cas d'usage, participation envisagée mais timing pas OK
- différents événements et sorties organisées pendant le stage 
  - TransIA : Tours, Présentations pluridisciplinaires sur l'intégration de l'IA (sous toutes ses formes) 
  dans la société
  - TAL-I : journée du GDR TaL-I à Paris (Jussieu) sur les travaux portant sur l'interprétabilité et 
  l'explicabilité dans le contexte des tâches de traitement de la langue naturelle
  - 2nd Symposium on Mental Health and AI : deux jours de conférence (suivi en pointillé en distanciel) 
  présentation des travaux à la croisée entre santé mentale et intelligence artificielle (impact, outils 
  de détection, de prise en charge)
  - Place du Numérique : événement public de sensibilisation et de communication sur le Numérique 
  organisé place du Martoi à Orléans (Des événements similaires organisés dans chaque département 
  de la région CVL) J'ai participé à la tenue du stand de l'université, été ambassadrice de l'événement 
  et ai joué le rôle d'un témoin dan sle tribunal des génération futures portant sur "L'IA nous 
  contrôle-t-elle?
  - Présentation aux journées d'équipe CA
    \end{verbatim}


\section{Outils utilisés durant le stage}

    \begin{verbatim}
- Développement de l'architecture et réalisation des expérimentations en python dans 3 formats 
(script, implémentation d'un CLI, notebook)
- Principales librairies python utilisées : PyTorch, TransformerLens, HuggingFace, PyMuPDF4LLM,...
- Utilisation de plusieurs git pour assurer le suivi des versions des outils développé lors du stage
- Pour l'execution de tâches demandantes en charge de calcul (Génération de la base de représentation 
des données, execution des différentes architectures explorées et/ou implémentées) : utilisation des 
grappes de calcul accessibles aux chercheurs du LIFO CaSciModOT et Mirev
- utilisation d'outils à base d'IA très modérée et pour les tâches de l'ordre de reformulation
(ChatGPT), traduction (DeepL), exploration de la littérature (LitMaps, Perplexity), rédaction 
de petites fonctions et outils python (ChatGPT)
    \end{verbatim}



\section{Contexte et enjeux du sujet}

    \begin{verbatim}
TODO        
    \end{verbatim}
